% SOURCE: STATE.md "Key data facts" and "Model"; Decision Log #1 (dedup
% before the split), #2 (macro-recall primary), #3 (RF/none hardcoded),
% #12 (the measured dedup penalty), #20 (adolescent ages retained),
% #21 (generic maternal-health risk only); related_work.md §5 (leakage
% wording: quantification, not discovery); results/tables/p1_model_metrics.csv
% and results/tables/p1_duplicate_leakage.csv.
% Every number below is traceable to one of those two committed tables or
% to a counted property of the cached raw CSV.
\section{Data and the Model Vehicle}
\label{sec:data}

\subsection{Dataset}
\label{sec:data-dataset}

We use the UCI Maternal Health Risk dataset (repository id 863), released
under CC BY 4.0. It contains 1014 records collected in rural Bangladesh,
each with six features --- age, systolic and diastolic blood pressure, blood
sugar, body temperature and heart rate --- and a three-level risk label:
low, mid, high. It is public data and no other data is used anywhere in this
work. No real patient record, de-identified or otherwise, appears in this
paper or in the repository that accompanies it.

The dataset supports \emph{generic maternal-health risk} and nothing
narrower. It carries no label for any specific condition, and none of the
laboratory, urinary or obstetric-history fields that would let one be
inferred. We therefore use the generic term throughout, in the
manuscript, in the code and in every rendered artifact, and no tier names a
condition. This matters for the third tier in particular: a spoken message
that named a diagnosis would be asserting something the data cannot support.

The label itself is worth stating plainly. It is a risk stratification
recorded alongside the readings, not a delivered outcome. Nothing in this
paper --- no figure, no rendered card, no spoken message --- should be read
as a prediction of what happens to any patient. What we demonstrate is an
explanation design over a model that has been fit to this public dataset.

\subsection{Deduplication, and why it happens before the split}
\label{sec:data-dedup}

Two preprocessing steps are applied, both before any train/test split, and
both chosen so that neither learns a statistic from the data.

First, two records report a heart rate of 7 beats per minute. These are
physiologically impossible and are dropped, leaving 1012 rows. Second, the
remaining set contains a large number of exactly repeated feature-and-label
vectors: 561 of the 1012 rows are exact duplicates of an earlier row. We
drop them, leaving \textbf{451 distinct records} --- 233 low, 106 mid, 112
high, an imbalance of about 2.2:1 between the largest and smallest class.

Both steps are fixed rules rather than fitted transformations, which is why
performing them before the split introduces no leakage: nothing about the
test fold informs what is dropped. Everything that \emph{does} learn from
data --- feature scaling, and SMOTE where a configuration used it --- is
fitted inside each cross-validation training fold only, through an
\texttt{imblearn} pipeline, never on the full dataset. We considered instead
retaining the duplicates and using group-aware cross-validation to keep
identical rows on the same side of every split. That is a valid alternative
and it preserves more rows; we rejected it because it adds machinery a
vehicle model does not need, and because deduplication makes the resulting
number directly comparable to what a reader would get from a plain
\texttt{drop\_duplicates}.

Deduplication does not resolve every ambiguity in the data, and we do not
claim it does. Among the 451 distinct records there remain \textbf{35
feature vectors that carry more than one risk label} --- identical six-value
readings recorded against different levels of risk. This is irreducible with
the features available: no model, however specified, can separate rows it
cannot distinguish. We retain these rows rather than adjudicating between
them, and we return to what they imply for the achievable ceiling in
Section~\ref{sec:limitations}.

\subsection{What the duplicates were worth}
\label{sec:data-leakage}

The duplicate rows are not merely untidy. Because a duplicated row can
appear in a training fold and in the test fold of the same split, retaining
them lets a model be scored partly on records it has already memorised. We
measured the size of that effect directly, by scoring the identical model
under the identical five-fold stratified cross-validation twice, changing
only whether duplicates were dropped (Table~\ref{tab:leakage};
\texttt{results/tables/p1\_duplicate\_leakage.csv}).

\begin{table}[t]
\centering
\caption{The same model and the same cross-validation, scored with exact
duplicate rows dropped and retained. Mean $\pm$ standard deviation across
five stratified folds.}
\label{tab:leakage}
\begin{tabular}{lrrrr}
\toprule
Duplicates & $n$ & Macro-recall & Accuracy & High-risk recall \\
\midrule
Dropped & 451  & $0.580 \pm 0.027$ & $0.641 \pm 0.025$ & $0.723 \pm 0.054$ \\
Kept    & 1012 & $0.859 \pm 0.026$ & $0.854 \pm 0.026$ & $0.908 \pm 0.031$ \\
\bottomrule
\end{tabular}
\end{table}

Retaining the duplicates moves accuracy from 0.641 to 0.854 and macro-recall
from 0.580 to 0.859. The inflated figure sits inside the accuracy band
typically reported for this dataset, and the reported range extends higher
still: Maisoon et al.\ report 83.74\% and Mamun et al.\ report
99.51\%~\citep{maisoon2026stratification,mamun2025identification}. Neither
paper describes duplicate handling; the published preprocessing on this
dataset is generally given as normalisation and the removal of
inconsistencies, and both report $n = 1014$.

We are careful about what this licenses us to say. Data leakage as a
phenomenon is well documented --- Kapoor and Narayanan set out its taxonomy
and its role in the reproducibility crisis across
fields~\citep{kapoor2023leakage}, Rosenblatt et al.\ show specifically that
leakage through repeated subjects inflates prediction performance and that
the effect is worst on small datasets, which is exactly our
situation~\citep{rosenblatt2024leakage}, and Bernett et al.\ give the
practitioner's checklist for avoiding it~\citep{bernett2024guiding}. The
contribution type is also established: re-analysing an inflated published
result to recover the true number is what Eltawil et al.\ and Young et al.\
do in their respective
domains~\citep{eltawil2026comment,young2025leakage}. \textbf{We do not claim
to be the first to notice these duplicates.} We searched for prior
documentation of them in this benchmark and found none, but absence of
evidence is weak here: a \texttt{drop\_duplicates} result is exactly the
kind of thing noticed in a notebook and never written up. What we claim is
the measurement. The published band on this dataset is recoverable only when
identical rows are allowed to straddle the split; after deduplication,
macro-recall on this vehicle is 0.580. We report the deduplicated figures
throughout and state the comparison here, up front, so that our lower
headline number is legible as hygiene rather than as a weaker model. It is
reproducible in either direction from the committed table.

\subsection{The model is a vehicle}
\label{sec:data-model}

We compared logistic regression, random forest and XGBoost, each with and
without in-fold SMOTE, under five-fold stratified cross-validation
(\texttt{results/tables/p1\_model\_metrics.csv}). The primary metric is
\textbf{macro-recall}, with macro-F1, per-class recall, one-versus-rest
ROC-AUC and the confusion matrix as support. Macro-recall is the right
primary here because the costs are asymmetric: failing to flag a high-risk
mother is the error that matters, and headline accuracy on an imbalanced
three-class problem conceals exactly that failure.

The final model is a \textbf{random forest with no resampling}
(\texttt{n\_estimators=300}, \texttt{random\_state=42}). It ties the
SMOTE-resampled random forest on macro-recall at 0.580, and we broke the tie
deliberately rather than by argmax: the unresampled model has tighter
variance across folds ($\pm 0.027$ versus $\pm 0.040$), the better
high-risk recall (0.723), and no dependence on a synthetic-oversampling
step. It also lets SHAP run on the six raw clinical features in their own
units, which is what makes the ASHA and mother tiers expressible as
templates over attributions at all. The selection is fixed as a constant in
the code, and the model-selection script still prints the argmax winner so
the tie is visible rather than papered over.

Macro-recall of 0.580 is the honest post-deduplication number on this
dataset, and we present it as the premise of the paper rather than as a
shortfall to be excused. The design question this work addresses --- how an
uncertain prediction should be rendered to readers of differing power ---
only exists because models on data like this are uncertain. A model that was
right 99\% of the time would not need the machinery we describe. The
attributions themselves are stable and unremarkable: blood sugar is the
leading global driver for every class, with systolic blood pressure second,
and diastolic pressure and heart rate close to noise
(\texttt{results/tables/shap\_global\_mean\_abs.csv}). We say ``blood sugar
and systolic pressure'' rather than ``the blood-pressure features''
throughout, because the data supports the first and not the second.

\subsection{Maternal age}
\label{sec:data-age}

One property of the dataset requires comment at first mention, because a
reader will otherwise meet it in a figure. The dataset contains very low
maternal ages --- 95 of the 451 distinct records are under 18, and the
minimum is 10 --- reflecting adolescent pregnancy in the source population.
We retain rather than filter these rows: excluding a group the system would
encounter is itself a harm. One consequence surfaces later in the paper. Our
designated failure case is a 13-year-old whose recorded body temperature of
101\,\textdegree F drives a true-low record to a predicted-high output, with
blood sugar and systolic pressure both arguing against that flip. That case
shows the model over-weighting a demographic feature, which is precisely why
final judgment stays with a human, and why the rendering rules of
Section~\ref{sec:tiers} exist. All data here is public (UCI Maternal Health
Risk); no real patient or child record is exposed.
